\documentclass[12pt]{article}
\usepackage[utf8]{inputenc}
\usepackage[T1]{fontenc}
\usepackage[margin=1in]{geometry}
\usepackage{setspace}
\usepackage{xcolor}
\usepackage[colorlinks=true, urlcolor=blue!60!black, linkcolor=blue!60!black]{hyperref}
\setlength{\parindent}{0em}
\setlength{\parskip}{0.45em}
\renewcommand{\baselinestretch}{1.0}
\begin{document}
\begin{center}
{\Large \textit{Katzenbach v. McClung}, 379 U.S. 294 (1964)\par}
\vspace{0.3em}
{\normalsize Case Brief --- One-Page Edition\par}
\end{center}

\textbf{Memory Jogger.} Congress may ban discrimination at Ollie's Barbecue, a local restaurant, because much of its food crossed state lines (\href{https://www.courtlistener.com/opinion/106951/katzenbach-v-mcclung/}{CourtListener}).

\textbf{Facts.} Ollie's Barbecue in Birmingham had refused dining-room service to African Americans since 1927. Forty-six percent of its food (\$69,683 of meat) came from outside Alabama. Title II of the Civil Rights Act of 1964 covered restaurants serving food a ``substantial portion'' of which had moved in commerce; no interstate travelers were shown.

\textbf{Procedural History.} A three-judge district court held Title II unconstitutional as applied. The Supreme Court reversed, alongside \textit{Heart of Atlanta Motel}.

\textbf{Votes.} 9--0: Clark for the Court; Black, Douglas, and Goldberg concurred.

\textbf{Holding.} Congress may regulate local commercial activity when it has a rational basis for finding that the class, in the aggregate, substantially affects interstate commerce; no formal findings or case-by-case proof is needed.

\textbf{Reasoning.} Segregation reduced interstate food purchases, depressed business, obstructed travel, and deterred new industry; under \textit{Wickard}, Ollie's small share, aggregated, was far from trivial. Black insisted the question is ``ultimately a judicial, rather than a legislative, question''; Douglas preferred \S~5 of the Fourteenth Amendment; Goldberg stressed ``the vindication of human dignity, and not mere economics.''

\textbf{Examples.} \emph{Same-Side:} local grocery stores with interstate inventory. \emph{Opposite-Side:} a knife ban near parks (\textit{Lopez}). \emph{Fence-Sitter:} a farm-to-table diner with only out-of-state salt.

\textbf{Critique.} Originalists see an unbounded commerce power; progressives regret resting equal dignity on meat shipments.

\textbf{Key Quotation.} ``Of course, the mere fact that Congress has said when particular activity shall be deemed to affect commerce does not preclude further examination by this Court. But where we find that the legislators, in light of the facts and testimony before them, have a rational basis for finding a chosen regulatory scheme necessary to the protection of commerce, our investigation is at an end.'' --- Justice Clark, 379 U.S. at 303--304.

\end{document}
